% Intermediate summary matrix for Q1. May later be condensed, moved to an appendix, or redrawn as a figure.
\begin{table}[tbp]
\centering
\footnotesize
\singlespacing
\caption{Working summary for Q1: what carries acoustic meaning. Status reflects the strength of evidence reviewed in this section, not a formal meta-analysis.}
\label{tab:summary-q1}
\begin{tabular}{@{}p{1.55in}p{3.7in}p{1.0in}@{}}
\toprule
\textbf{Claim} & \textbf{Finding} & \textbf{Status} \\
\midrule
Natural $>$ synthetic & Often true vs noise or poorly matched controls; fails when frequency or timing is matched & Mixed \\
\addlinespace[2pt]
Naturalness is causal & Frequency band and gross timing often suffice; modulation adds to, rather than creates, responses & Strong against \\
\addlinespace[2pt]
Significance $\times$ intensity & Natural-call advantage can appear at high SPL and vanish at lower SPL & Thin (1 study) \\
\addlinespace[2pt]
Synthetic can acquire meaning & Maternal experience retunes responses to pup-call interval morphs & Strong (maternal) \\
\addlinespace[2pt]
Field call structure & Urgency (marmot rate, nonlinear noise); predator-specific escape (prairie dogs); call type and syllable position (ground squirrels) & Strong \\
\addlinespace[2pt]
Source ranking & Indirect social cues can outweigh predator vocalizations & Caution \\
\addlinespace[2pt]
Adequate control & Match level, spectrum, envelope, and modulation; test $>$1 SPL; know receiver experience & Gap in field \\
\bottomrule
\end{tabular}
\end{table}
